\documentclass[11pt,a4paper]{article}
\usepackage[T1]{fontenc}
\usepackage[utf8]{inputenc}
\usepackage{amsmath,amssymb,amsthm}
\usepackage[margin=1in]{geometry}
\usepackage[colorlinks=true,linkcolor=blue,urlcolor=blue]{hyperref}
\theoremstyle{plain}
\newtheorem{theorem}{Theorem}
\newtheorem{lemma}[theorem]{Lemma}
\newtheorem{proposition}[theorem]{Proposition}
\newtheorem{corollary}[theorem]{Corollary}
\theoremstyle{definition}
\newtheorem{remark}[theorem]{Remark}

\title{Recovering the unwritten recurrences of the table OEIS A237150}
\author{Adrian Perez Fontelles\\ \small Independent researcher}
\date{2 September 2026}
\begin{document}
\maketitle

\begin{abstract}
OEIS A237150 is a two-dimensional table $T(n,k)$ whose empirical recurrences are, for several of
its lines, given only as an ORDER --- ``k=2: [order 70]'' and the like --- with the
coefficients in a linked file rather than in the entry. Those conjectures can still be settled,
and the recurrences recovered. A line of the table is a fixed-width or fixed-height array
count, hence a walk count on finitely many vertices, hence satisfies some monic recurrence of
order at most that number; Berlekamp--Massey on twice as many exact terms returns the MINIMAL
one. Where its order equals the order the entry states --- and where the same holds after the
first few terms are dropped --- any recurrence of that order the line satisfies has a
characteristic polynomial that is a multiple of the minimal one and of equal degree, hence is
that polynomial. This note settles column 2.
\end{abstract}

\noindent\small 2020 Mathematics Subject Classification. 05A15, 11B37, 15A18.\normalsize

\section{The table and the unwritten conjectures}

OEIS A237150 is ``T(n,k)=Number of (n+1)X(k+1) 0..3 arrays with the maximum plus the lower median of every 2X2 subblock equal.''. It is read by antidiagonals and begins
\[
256, 1472, 1472, 8128, 17616, 8128, 48672, 197724,\ \dots
\]
The entry carries, among its empirical claims:
\begin{quote}\small
k=2: [order 70]
\end{quote}
As of the ``Last modified'' line on the live entry (Jul 23 2025 09:53:45, revision 6) these are still
recorded as empirical, and the recurrences themselves are not in the entry text.

\section{Why a line of the table is a walk count}

Fix the index that the line fixes. The entry defines $T(n,k)$ as a number of arrays of a shape
determined by $n$ and $k$, subject to a condition local to a bounded window of consecutive
lines. Substituting the fixed index into the entry's own wording turns the definition into an
ordinary one-parameter array count, and for such a count the admissible windows are the
vertices of a finite digraph and an array is a walk. So the line is a walk count on $S$
vertices, and by the Cayley--Hamilton theorem it satisfies a monic integer recurrence of order
at most $S$.

\section{Recovering the recurrences}

\begin{lemma}\label{lem:bm}
A sequence known to satisfy some linear recurrence of order at most $S$ is determined, with its
minimal recurrence, by its first $2S$ terms; Berlekamp--Massey applied to them returns that
minimal recurrence exactly.
\end{lemma}

\subsection*{Column $k=2$}
Substituting the fixed index into the entry's wording gives an ordinary one-parameter array
count, a walk on $S=216$ vertices. Berlekamp--Massey on $2S=432$ exact terms
returns a minimal recurrence of order $70$ --- the order the entry states --- with
integer coefficients
\[
(c_1,\dots,c_{70})=(39,\ -160,\ -12044,\ 145902,\ 1301357,\ -28998486,\ -23143247,\ 2983103077,\ -8372689257,\ -186522916060,\ 1021940846958,\ 7344299435320,\ -63749321988945,\ -166330077893225,\ 2599781261562279,\ 703737238969351,\ -74689355357780259,\ 96224930618884890,\ 1554555314427677603,\ -4037104924465079240,\ -23439383477887989510,\ 94468254131551344798,\ 245986355248174699197,\ -1538068456145135501031,\ -1509409575078360285981,\ 18585206013649610854908,\ -969679345185840587997,\ -171019089996424356270276,\ 141770311292210122719180,\ 1208287224153788793868230,\ -1833180013932432304366755,\ -6515970743144530671621596,\ 14722519233845227915335344,\ 26100016339490264593135930,\ -85911886883585776474149659,\ -71381602952343952813346574,\ 382888187578277599439919469,\ 88741065470822620977868868,\ -1330364967974923022020690482,\ 263612103359188023210204998,\ 3632822012287157342790093062,\ -1959772874708864128893661132,\ -7800203620907892204716451725,\ 6553800746731419550738384340,\ 13090078467385195437896000593,\ -14918802129637674735916961289,\ -16929578769554397337783200256,\ 25163812912874646295081404754,\ 16411352874801018496786413612,\ -32344378109509927413894038376,\ -11228174964373034514526126784,\ 31947787660331753005976517728,\ 4525159837076488065598251776,\ -24213200715632083926698812032,\ 13921824249548065145684992,\ 13966608945681628925157484544,\ -1371967094177676433883704320,\ -6046924385419647588041842688,\ 976284382984348699783348224,\ 1926631847549153169292263424,\ -376832937576072420213325824,\ -439760960432071967326076928,\ 89346900076942136064344064,\ 69296632458273967331868672,\ -12830845605765831928578048,\ -7130511556499570128060416,\ 1015785531774837102477312,\ 432510443602795258970112,\ -33789326924988071018496,\ -11836509492045108215808).
\]
so that $a(m)=\sum_{i=1}^{70}c_i\,a(m-i)$ for every $m$. Removing the first
$70+4$ terms and repeating gives order $70$ again, which is what the
identification below needs.

\begin{theorem}
For each line above, the displayed recurrence holds for every index, and it is the recurrence
the entry records.
\end{theorem}

\begin{proof}
It holds by Lemma~\ref{lem:bm}. For the identification, the entry asserts that the line
satisfies SOME recurrence of the stated order, possibly only from some index on. Let $q$ be its
characteristic polynomial and $q_{\min}$ the minimal polynomial of the tail. Then
$q_{\min}\mid q$, and the second Berlekamp--Massey run --- on the line with its first few
terms removed --- shows $\deg q_{\min}$ equals the stated order, which is $\deg q$. Both are
monic, so $q=q_{\min}$.
\end{proof}

\section{Verification}

Three checks.

First, each line's model was compared against the entry's own published values for that line,
taken from the antidiagonal DATA read in either orientation and from the ``Table starts''
block, and agrees at every available term. This is the only point at which reading the entry's
English enters, and a misreading gives different counts.

Second, each recovered recurrence was evaluated directly on the model's values, with no
Berlekamp--Massey involved, and holds at every index tested.

Third, each order was computed twice by independent routes --- modulo a $61$-bit prime, where
every intermediate is a machine word, and again in exact rational arithmetic --- and once more
on the line with its first few terms dropped, which is what the identification requires. All
agree.

\begin{thebibliography}{9}
\bibitem{oeis} The OEIS Foundation, \emph{The On-Line Encyclopedia of Integer
Sequences}, \url{https://oeis.org}, 2026. Sequence A237150.
\bibitem{massey} J.~L.~Massey, Shift-register synthesis and BCH decoding, \emph{IEEE
Trans. Inform. Theory} \textbf{15} (1969), 122--127.
\bibitem{stanley} R.~P.~Stanley, \emph{Enumerative Combinatorics}, Volume 1, 2nd ed.,
Cambridge University Press, 2011. (Transfer-matrix method, Section 4.7.)
\bibitem{horn} R.~A.~Horn and C.~R.~Johnson, \emph{Matrix Analysis}, 2nd ed., Cambridge
University Press, 2013.
\end{thebibliography}

\end{document}
